% !TeX root = ../教材题库_全部条目_转写工作稿.tex
% Source: XJ-B1-C04-EX-P097-Y0403
% 第4章 / 4.1 实数指数幂和幂函数 / 教材 p.97 / PDF 第104页
% OCR未校对：公式与文字均可能有误；未完成数学TeX转写；禁止直接发给学生
% 原题及相邻原解截图（核对用）：../题目原图/XJ-B1-C04-EX-P097-Y0403-01.webp
% 本文件只在转写完成后作为题干使用；OCR 条目尚不是可靠题干。
\begin{bankitem}{4.1 例题 例题5}{第4章 / 4.1 实数指数幂和幂函数 / 教材 p.97 / PDF 第104页}{XJ-B1-C04-EX-P097-Y0403}{OCR未校对：公式与文字均可能有误；未完成数学TeX转写；禁止直接发给学生}
\begin{ocrdraft}
[教材p.97 / PDF第104页]
\par
00用分数指数幂的形式表示下列根式的化简结果（式中字母都是正数）：
\par
3
\par
2
\par
（3），2（，）
\par
3
\par
解（1)
\par
36\ocrglyph{00AF}a26．一6
\par
磊3
\par
（3）卫2（3'）3一()y，2()y，广厂一00厂一x了3骁
\end{ocrdraft}
\end{bankitem}
